% =====================================================================
%  Your Agent is Underperforming: The Case for Agent Performance
%  Improvement Plans  ---  6-page ICA 2026 version
%
%  Cuts applied on top of the 11-page revision:
%    - VI-B1 pilot vignettes compressed to one paragraph (judge/baseline
%      points folded into Limitations)
%    - Table III + Section V replaced by one schema paragraph inside
%      Section IV (full schema -> supplementary material)
%    - Section II-E merged into the adversarial-subversion paragraph
%    - Tier-mismatch / mesh-disruption duplication pass
%    - Abstract <=150 words, 4 tight contribution bullets,
%      conclusion <=120 words
%    - Section III mode paragraphs compressed (Table I is canonical)
%    - Table II -> supplementary material (pointer sentence kept)
%    - II-A / II-B / IV-A / IV-D paragraph-level tightening; Fig. 1 cut
%    - References 34 -> 28
% =====================================================================

\documentclass[conference]{IEEEtran}
\IEEEoverridecommandlockouts

\usepackage{cite}
\usepackage{amsmath,amssymb,amsfonts}
\usepackage{graphicx}
\usepackage{textcomp}
\usepackage{xcolor}
\usepackage{booktabs}
\usepackage{array}
\usepackage{url}
\usepackage[htt]{hyphenat}   % lets long \texttt identifiers break across lines

% Placeholder macro: every unresolved result is wrapped in \ph{} so the
% remaining TODOs are greppable and visibly red in the draft.
% Before submission: \renewcommand{\ph}[1]{#1} after filling numbers in.
\newcommand{\ph}[1]{\textcolor{red}{[#1]}}
\newcommand{\code}[1]{\texttt{\small #1}}

\begin{document}

\title{Your Agent is Underperforming:\\
The Case for Agent Performance Improvement Plans}

\author{\IEEEauthorblockN{Anonymous Submission}
\IEEEauthorblockA{\textit{Author information withheld per review guidelines}}}

\maketitle

% =====================================================================
\begin{abstract}
Production AI agents frequently degrade after deployment, yet structured
remediation culture is largely absent. We introduce the Agent Performance
Improvement Plan (APIP): a lifecycle protocol adapted from organizational
management that governs the remediation of underperforming agents through
statistically disciplined triggers on a behavioral contract, measurable cause
attribution over a five-category failure-mode taxonomy, tiered interventions, a
bounded window with formal exit criteria, and a declarative record meeting
governance traceability requirements. We evaluate it end-to-end in two live
multi-agent environments---customer service ($\tau^2$-bench) and software
engineering (ChatDev)---injecting failure modes with known ground truth and
scoring against no-action, always-Tier-1, and oracle baselines. Under a
stale-corpus drift injection, APIP and the oracle alone exit durably: the
always-Tier-1 responder patches four times, regresses three times, and ends
indistinguishable from no action---evidence that time-to-resolution without a
durability criterion rewards the worst strategy. Attribution is correct in all
three non-abstaining cycles. We further report a null result: inserting a
competent specialist into a role-differentiated mesh produced no peer
degradation, localizing the hypothesized crowding mechanism to architectures
with a bounded shared context.
\end{abstract}

\begin{IEEEkeywords}
AI agents, post-deployment monitoring, remediation, behavioral drift, AI
governance, multi-agent systems, performance improvement plan
\end{IEEEkeywords}

% =====================================================================
\section{Introduction}

Organizations increasingly deploy conversational agents as their main customer
interfaces, handling support tickets, returns, and billing disputes at scale.
Pre-release evaluation is well-equipped---standardized benchmarks
\cite{liang2022}, red-teaming \cite{ganguli2022}, RLHF alignment
\cite{ouyang2022}---but there is a notable lack of culture around
post-deployment remediation: a comprehensive study of deployed agents found
many fail or underperform, with limited understanding of how organizations
react \cite{pan2025}.

Ad-hoc responses to degraded agents are standard: engineers ``fix'' individual
prompts, product managers file tickets, systems get rolled back---with no
structure, no documentation, and no evaluation of whether the fixes worked.
Agents recover through informal means or are quietly removed. Beyond the
operational cost, this is a governance risk: the EU AI Act requires
documentation and human oversight of deployed AI systems \cite{euaiact}, and
Wachter et al.\ \cite{wachter2017} note that accountability in automated
decision-making requires traceability, which formalized remediation can provide.

We recommend the performance improvement plan (PIP) as the model: an
established HR tool containing the elements common to accountability
protocols---a trigger, cause identification, a bounded corrective window,
check-ins, and an explicit definition of recovery or exit. Kirkpatrick's
\cite{kirkpatrick1994} four-level model---reaction, learning, behavior,
results---complements it by evaluating whether remediation produced lasting
change.

We make the following contributions:

\begin{itemize}
  \item We introduce the APIP, a formalized post-deployment remediation
  lifecycle protocol grounded in organizational management and training
  evaluation research.

  \item We develop a five-category failure-mode taxonomy with a published
  deterministic mapping from trace-level annotations (MAST \cite{cemri2025}) and
  disaggregation signatures onto taxonomy categories and intervention tiers.

  \item We propose a three-tier intervention hierarchy with escalation
  criteria, evaluated through Kirkpatrick's framework, and a declarative APIP
  schema for tooling and audit trails.

  \item We evaluate the full protocol in two live multi-agent environments
  against a four-condition baseline family, report the first attribution
  accounting for a remediation protocol, and localize the mesh disruption
  problem to architectures with a bounded shared context via a null result.
\end{itemize}

% =====================================================================
\section{Background and Related Work}

\subsection{Agent Evaluation and the Deployment Gap}

Pre-deployment evaluation has matured---benchmarks \cite{liang2022},
red-teaming \cite{ganguli2022}, RLHF \cite{ouyang2022}---while post-deployment
monitoring has received far less attention. Pan et al.\ \cite{pan2025} find in
the first large-scale study of production agents that most deployments fail or
underdeliver, with little known about post-release management; AgentCompass
\cite{sapra2025} detects but prescribes no remediation; Rath \cite{rath2026}
projects a 42\% task-success reduction under unchecked drift. On the diagnosis
side, MAST \cite{cemri2025} supplies an empirically derived failure taxonomy
with an LLM-annotator pipeline; it classifies \emph{what went wrong in a trace}
without selecting a remediation, and we adopt it as the trace-labeling stage of
cause attribution (Section~\ref{sec:attribution}). Run-level toolkits such as
AgentDebugX \cite{zhu2026} close a detection--attribution--recovery--rerun loop
for a single failed execution; the APIP operates one level above, deciding
whether an aggregate degradation warrants intervention, at what tier, and until
when---a successful rerun being Level~1/2 evidence
(Section~\ref{sec:kirkpatrick}), not durable recovery against the contract.

\subsection{The Performance Improvement Plan as a Structural Template}

The PIP is a formal document specifying observed deficiencies, root causes,
required corrective actions, a time window, a check-in schedule, and success
criteria \cite{kirkpatrick2016}. The analogy needs one qualification: a human
employee can self-direct improvement, a deployed agent cannot, so the APIP is
imposed entirely by operators---the agent is the subject, not the author, of
the plan---collapsing the manager--employee dyad into a single operator who
both writes and implements it.

\subsection{Kirkpatrick's Four-Level Evaluation Framework}
\label{sec:kirkpatrick}

Kirkpatrick's four-level model \cite{kirkpatrick1994, kirkpatrick2016}
evaluates interventions hierarchically---Reaction, Learning, Behavior
(transfer to on-the-job change), Results (organizational outcomes)---and we
adopt it to ask of each intervention: did it change immediate output patterns
(Level 1/2), produce durable change on the failure-mode class (Level 3), and
improve the contract's business metrics (Level 4)? An intervention improving
Level 1/2 without Level 3/4 has not remediated the agent; it has suppressed the
symptom. Our always-Tier-1 baseline is exactly this failure, measured.

\subsection{AI Governance and Documentation Requirements}
\label{sec:governance}

Regulation increasingly requires documentation, oversight, and traceability:
the EU AI Act \cite{euaiact} mandates post-market monitoring and logs
sufficient to identify sources of error, and Kolt et al.\ \cite{kolt2024} argue
accountability requires linking actions to specific agents and operators
through identifiers, monitoring, and activity logging. The schema's
\code{apip\_id}, \code{agent\_id}, \code{owner}, and \code{exit\_evaluated\_at}
fields operationalize that linkage, positioning the APIP as a governance
instrument.

% =====================================================================
\section{A Taxonomy of Agent Failure Modes}
\label{sec:taxonomy}

\begin{table}[t]
\caption{Agent failure mode taxonomy with Kirkpatrick level mapping and
remediation tier.}
\label{tab:taxonomy}
\centering
\scriptsize
\begin{tabular}{@{}p{1.4cm}p{4.1cm}p{1.7cm}@{}}
\toprule
\textbf{Failure Mode} & \textbf{Description \& Signals} & \textbf{Remediation
Tier} \\
\midrule
Behavioral Drift &
Gradual degradation as world-state diverges from training or retrieval context.
Signals: rising hallucination rate on new entities, outdated policy citations.
Kirkpatrick Level 4 (Results). &
Tier 2: Retrieval / RAG update \\
\addlinespace
Input Brittleness &
Strong on modal inputs, catastrophic on a specific subclass. Signals: bimodal
CSAT, failure clustering by segment or topic. Level 2 (Learning/Capability). &
Tier 1 (examples) or Tier 3 (fine-tune) \\
\addlinespace
Alignment Creep &
Subtle drift in tone, value expression, or policy compliance. Signals: rising
review flags, brand guideline violations. Level 3 (Behavior). &
Tier 1: Prompt constraint injection \\
\addlinespace
Tool Misuse &
Correct intent, incorrect tool invocation. Signals: action errors, silent
downstream failures. Level 3 (Behavior). &
Tier 1 (prompt) or Tier 3 (fine-tune) \\
\addlinespace
Exogenous Disruption &
Peer agents degrade after a mesh change event; no intrinsic mode found in them.
Signals: degradation concentrated in peer roles, onset at a mesh-composition
event. Level 4 (Results), system level. &
Mesh protocol (Sec.~\ref{sec:mesh}), not an agent-level tier \\
\bottomrule
\end{tabular}
\end{table}

A critical prerequisite to structured remediation is structured cause-finding.
We propose a five-category taxonomy for customer-facing agents
(Table~\ref{tab:taxonomy}): four intrinsic modes---behavioral drift, input
brittleness, alignment creep, tool misuse---plus one exogenous mode, mesh
disruption (\code{exogenous\_disruption}), in which the failing agent is not
the cause of its own degradation. The categories suggest distinct remediation
pathways---the operationally important property---each corresponding to a
different level of the agent's behavioral stack and consequently a different
intervention tier.

Two properties deserve emphasis. Input brittleness is masked by aggregate
metrics---an agent at 85\% overall success may fail catastrophically on the 8\%
of requests from one segment---so the trigger must monitor disaggregated
splits. And behavioral drift is a freshness failure that retrieval refresh
remediates and fine-tuning does not, whereas alignment creep is value
misalignment, distinct from capability failure \cite{hendrycks2021}.

\textbf{Exogenous Disruption} is the taxonomy's escape hatch: attribution
correctly concludes that \emph{no} intrinsic mode explains the degradation,
the cause lying in a mesh change---most prominently a new, individually
competent peer. Its signature is structural rather than behavioral, and its
remediation is a mesh-level protocol rather than an agent-level tier
(Section~\ref{sec:mesh}).

\subsection{From Trace Labels to Taxonomy Categories}
\label{sec:mapping}

The taxonomy is operational only if a degradation window can be classified by
procedure rather than intuition. Attribution combines a trace-level label
distribution from MAST \cite{cemri2025} with a \emph{disaggregation
signature}---how degradation distributes across task categories, roles, and
time---through a published deterministic mapping (supplementary material); two
mapping rows share a MAST cluster but differ on signature, which is why the
taxonomy is two-level. The mapping is frozen before any evaluated run, and
below-confidence combinations \code{abstain}, logged and scored as attribution
failures rather than silently resolved.

A further category, \textbf{Adversarial Subversion}---failures whose proximate
cause is an adaptive attack rather than capability or value drift
\cite{bhatt2025, schoen2025}---we declare out of scope: none of the three tiers
addresses an adversarial dynamic, and the effective responses belong to the AI
Control literature \cite{greenblatt2023}, which decides per action under a
threat model rather than over an aggregated contract window. Control audits can
feed APIP triggers, and the schema reserves the
\code{adversarial\_subversion} value.

% =====================================================================
\section{The APIP Framework}
\label{sec:framework}

We now describe the APIP as a structured lifecycle protocol of five phases:
Trigger, Cause Attribution, Intervention Planning, Remediation Window, and Exit
Evaluation.

\subsection{Phase 1: Trigger}
\label{sec:trigger}

An APIP initiates when performance crosses a trigger threshold on monitored
metrics. This presupposes a \emph{behavioral contract}---a formal
specification, written at deployment time, of task surface, success metrics and
measurement methodology, acceptable ranges, escalation policy, and human review
points; most teams currently skip it, delegating expectations to informal
organizational memory. Two trigger modalities are both necessary: absolute
floors catch severe degradation, while relative drift triggers catch gradual
decline that may never breach a floor.

\subsubsection{Statistical discipline of the trigger}
Implementing the trigger phase against live agents exposed a failure mode of the
naive design that we consider a finding in its own right. Thresholds must be
\emph{derived}, not hand-picked: each contract metric's envelope is estimated
from a calibration window as mean $\pm\,k\sigma$ (with $k$ reported), which
preempts the objection that thresholds were chosen so that triggers would fire
where desired. Even then, a single-gate rule---a rolling mean against the derived floor, across
the 15--25 metrics and disaggregated splits a realistic contract
monitors---breached in 75\% of \emph{clean, zero-injection} runs, because a small
calibration sample is treated as exact and the per-tick tests are uncorrected. A
breach must therefore pass two gates: the degradation must be at least the
contract's $k\sigma$ (it matters), and it must survive Benjamini--Hochberg
adjustment across the metric family (it is real). Disaggregated groups
additionally require a minimum per-group episode count; under-powered groups
receive no envelope, and the exclusion is recorded rather than silent. The
detector's significance level is swept on clean-arm runs only and frozen before
any evaluated window---clean data carries no information about where triggers
\emph{should} fire, so this does not compromise pre-registration. So calibrated,
the detector raised no false alarm on any clean tick we measured
(Section~\ref{sec:evaluation}). A trigger phase without this discipline
initiates APIPs on noise, and every downstream phase inherits the error.

\subsection{Phase 2: Cause Attribution}
\label{sec:attribution}

Once triggered, the APIP requires a structured cause attribution process before
any intervention is applied. The goal is to identify which failure mode category
(Section~\ref{sec:taxonomy}) best explains the degradation. Tier mismatch from
misdiagnosis is the primary source of avoidable regression: prompt engineering
applied to a retrieval freshness problem, for instance, consumes engineering
cycles with no improvement.

This is distinct from, and composes with, automated failure attribution, which
localizes the responsible agent or step within a single failed trajectory
\cite{ma2025attribution, zhang2025, agenttrace2026}. APIP Phase 2 operates one
level of aggregation above, classifying a window-level degradation into a
remediation-oriented category to select a tier; trajectory-level outputs are
natural evidence inputs, but localizing the failing step falls short of
diagnosing the underlying cause---the gap the mandatory taxonomy classification
closes at the deployment level.

The process should include disaggregated performance analysis to isolate
failure locus; human-adjudicated analysis of sampled failure cases, applying
automated trajectory-level attribution where tooling permits \cite{zhu2026,
ma2025attribution, agenttrace2026}; retrieval audit if drift is suspected; tool
invocation log analysis if misuse is suspected; and structured classification
into the taxonomy, closing with a written cause attribution report that becomes
part of the APIP record.

\subsubsection{A measurable instantiation: the three-stage funnel}
\label{sec:funnel}
We instantiate this as a three-stage funnel scoreable against ground truth:
(1)~every contract metric is split by task category, agent role, and user
segment, producing a disaggregation signature; (2)~a MAST annotator
\cite{cemri2025} labels a sample of failing traces, with the judge required to
differ from the agent's backbone---preferably by model family---to avoid a
judge sharing the failure modes it must detect; (3)~the frozen mapping
(Section~\ref{sec:mapping}) combines signature and label distribution into a
\code{failure\_mode} with a confidence score, returning \code{abstain} below
threshold. The mapping table is hashed into every run manifest, making post-hoc
edits detectable, and the funnel's discrete output is scored as a confusion
matrix against injected ground truth (Section~\ref{sec:evaluation}), with
abstention rate reported alongside---never instead of---accuracy.

\subsection{Phase 3: Intervention Planning and the Tier Model}

Interventions are organized into three tiers of increasing cost, disruption,
and lead time; cause attribution selects the starting tier, with escalation
permitted within the window. \textbf{Tier 1 (prompt-level)}: constraint
injection, persona adjustment, example replacement---fastest, least disruptive,
easiest to roll back; the first response to alignment creep and much input
brittleness. \textbf{Tier 2 (retrieval and context)}: RAG index updates,
knowledge-base revisions, tool schema fixes---addressing drift from world-state
divergence and tool misuse from stale descriptions, with contract validation
before re-deployment. \textbf{Tier 3 (model-level)}: fine-tuning, DPO, or
replacement---highest cost and regression risk, warranted when lower tiers are
exhausted or the failure mode is embedded in model weights.

\subsection{Phase 4: Remediation Window and Check-In Cadence}

The remediation window bounds the period during which interventions are
implemented and evaluated, creating urgency and a structured decision point.
Duration is calibrated to the tier---7--14 days for Tier 1, 30--60 for Tier
3---mirroring organizational PIP durations of 30--90 days \cite{kirkpatrick2016}.

Check-ins review performance against the recovery trajectory at multiple
Kirkpatrick levels---flagging interventions that improve surface metrics
without Level 3/4 change as potentially superficial---and act as early
warnings: no improvement at the first check-in triggers tier escalation rather
than waiting for expiry. A recommended cadence is 25\%, 50\%, and 75\% of
elapsed time.

\subsection{Phase 5: Exit Evaluation}
\label{sec:exit}

At window close, a formal exit evaluation records one of five outcomes:
\emph{resolved} (all contract metrics returned to envelope); \emph{regressed}
(envelope restored but breached again---the signature of a tier-mismatched
intervention that suppressed symptoms); \emph{timeout} (window expired without durable
recovery; more fundamental intervention warranted); \emph{retired} (the agent
cannot be returned to contract at acceptable cost); and \emph{recalibrated}
(attribution concluded the \emph{environment} legitimately changed---a moved
goalpost---so the contract is re-baselined rather than the agent treated). We
deliberately do not score recalibration as a resolution: conflating the two
would let re-baselining masquerade as remediation success. Window
\emph{extension} is an in-protocol check-in decision, not an exit outcome, so
every APIP terminates in exactly one of the five states; the record becomes
part of the organizational AI audit trail.

\subsection{The APIP Record as a Declarative Schema}
\label{sec:schema}

Each APIP is recorded as an instance of a declarative schema (full field
specification in supplementary material). Four field groups do architectural
work: \code{behavioral\_contract\_ref} enforces writing a contract at
deployment time, without which trigger and exit cannot be operationalized;
\code{kirkpatrick\_levels} captures multi-level check-in data, enabling
post-hoc analysis of superficial versus durable improvement;
\code{mapping\_table\_hash} and \code{environment\_provenance} extend the
visibility mechanisms of Kolt et al.\ \cite{kolt2024} to \emph{decision-rule}
linkage---an auditor can verify the attribution rules and environment versions
actually in force, converting pre-registration from assertion to checkable
property; and \code{mesh\_context} (Section~\ref{sec:mesh}) supports
multi-agent cause attribution.

% =====================================================================
\section{Evaluation in Two Live Multi-Agent Environments}
\label{sec:evaluation}

We evaluate the APIP end-to-end---calibration, trigger, attribution, tier
selection, bounded window, exit---in two live multi-agent environments with
real LLM inference and controlled failure-mode injection against known ground
truth. $\tau^2$-bench (\code{banking}) \cite{barres2025} matches the
customer-service domain motivating this work, scores every action against
deterministic database state, and has a retrieval corpus that can be staled
independently of task ground truth. ChatDev \cite{qian2024} supplies a
role-differentiated mesh (CEO $\rightarrow$ CTO $\rightarrow$ Programmer
$\rightarrow$ Reviewer $\rightarrow$ Tester), the substrate for creep,
disaggregated brittleness detection, and mesh insertion. That one
contract/trigger/tier/exit machinery governs two unrelated domains is the
generality argument for the contract abstraction.

\subsection{Shared Harness and Experimental Design}


Both environments plug into a common controller through an adapter interface;
traces are normalized to a shared span format, so the attribution pipeline is
written once and reused unmodified. Each run proceeds in ticks: a calibration
window derives every contract threshold as mean $\pm\,k\sigma$
(Section~\ref{sec:trigger}); failure modes are injected on a pre-registered
schedule; both trigger modalities evaluate at every tick; breaches pass through
the attribution funnel; and the selected intervention runs in a bounded window
with check-ins at 25/50/75\%, closing in one of the five exit outcomes.

\textbf{Baseline family.} The drift cell runs under four conditions:
\emph{APIP} (the full protocol); \emph{B1 no-action} (degradation runs to
horizon); \emph{B2 always-Tier-1} (a scripted prompt patch regardless of
failure mode---the ad-hoc pattern typical of production practice); and
\emph{B4 oracle} (ground-truth mode revealed at trigger, correct tier applied
immediately). B1 and B4 bracket the achievable range; B2 is the comparison that
matters, because it is what organizations actually do. Injected ground truth is
quarantined from the controller---written into the run record for scoring and
revealed only to B4---and a test enforces that the APIP condition never
branches on it.

\textbf{Statistical discipline.} Injection schedules, post-calibration
thresholds, the frozen mapping table, and the analysis plan were fixed before
the evaluated runs, and judge models are independent of backbone models by
construction. Each reported condition is a single seed, so we report observed
trajectories rather than confidence intervals and treat replication as the
principal limitation, rather than claim a variance estimate the runs do not
support.

\subsection{$\tau^2$-bench: Behavioral Drift with Deterministic Ground Truth}
\label{sec:taubench}

The primary drift experiment injects nothing into the agent; instead the
\emph{world} moves. Bank policy clauses and task ground truth advance on a
deepening schedule (one additional clause every five ticks from tick 12) while
the agent's retrieval corpus stays frozen---operationalizing the most familiar
production drift incident, an agent confidently citing superseded policy
because its knowledge base was never re-indexed, with the deterministic
database scoring every action objectively.

\begin{table}[t]
\caption{$\tau^2$-bench drift cell, all four conditions on one injection
schedule and task stream. TTR is reported only where a durable exit was
achieved. B2's four cycles each closed quickly; none held.}
\label{tab:baselines}
\centering
\footnotesize
\begin{tabular}{@{}lcccccl@{}}
\toprule
\textbf{Cond.} & \textbf{Det.} & \textbf{Act} & \textbf{Exit} & \textbf{TTR} &
\textbf{Cyc.} & \textbf{Outcome} \\
\midrule
B4 oracle      & 17 & 17   & 19 & 2  & 1 & resolved \\
APIP           & 17 & 27   & 29 & 12 & 1 & resolved \\
B2 Tier-1      & 17 & 17   & 39 & -- & 4 & timeout \\
B1 no-action   & 17 & --   & 39 & -- & 1 & timeout \\
\bottomrule
\end{tabular}
\end{table}

Detection fired at tick 17 in every condition---a lag of five ticks on the
judge-scored policy citation error rate, reproduced across three seeds---and
the calibrated detector raised no false alarm across twenty clean ticks at any
swept operating point. Under APIP, attribution returned
\code{behavioral\_drift} at confidence 0.65 and the matched Tier 2 corpus
refresh resolved the cycle in 12 ticks against the oracle's 2. That gap is
diagnosis, but ten of its twelve ticks are the re-attribution cadence---a
tunable protocol parameter---rather than diagnostic difficulty, since the
diagnosis was correct at first attempt. A recurring pilot failure sharpened the
exit test: a Tier 2 refresh that was real but partial closed the ticket without
fully rebuilding the index, and the contract re-breached within four ticks.
Because exit evaluation re-measures the envelope rather than trusting the
intervention log, the protocol verifies that remediation was material and
complete, not merely performed.

B2 reproduced the ad-hoc pattern the protocol is meant to displace. Its Tier 1
prompt patch names the clauses amended at patch time, recovers, and re-breaches
as the ladder deepens---four cycles, three regressions, and a final
\code{timeout} that leaves it indistinguishable from taking no action at all.
APIP and the oracle were the only conditions to exit durably.

This exposes a measurement trap. Scored on per-cycle time-to-resolution alone
B2 outperforms APIP, because each shallow patch closed fast; it also never
fixed the problem. Any evaluation of remediation reporting mean TTR without a
durability criterion will systematically favor the worst strategy, so we treat
\emph{durable exit} as the outcome measure and report TTR only where one was
achieved.
\subsection{ChatDev: Alignment Creep and Input Brittleness}
\label{sec:chatdev}

\textbf{Alignment creep} perturbs the Reviewer persona cumulatively (``keep
reviews brief'' $\rightarrow$ ``avoid blocking progress'' $\rightarrow$
``approve unless catastrophic''), one rung every three ticks. This is the taxonomy's
prediction under test: build success never left 1.00, so no absolute floor was
crossed and detection was available only through the drift modality, which
fired at the injection tick itself (lag 0) on review engagement. Attribution
abstained while a single rung was visible, then returned \code{alignment\_creep}
at confidence 0.73 once the pattern accumulated, selecting a Tier 1 charter
re-injection; review engagement recovered from 6.85 to 7.76 within one tick and
held twelve, exiting \code{resolved} at a TTR of 17. A companion cell attributed
the same mode correctly at 0.61 but never cleared its exit envelope, exiting
\code{timeout}: correct diagnosis is necessary, not sufficient.

\textbf{Input brittleness} strips the Programmer's GUI exemplars from its role
prompt, scoped to GUI-category episodes. Targeted requirement coverage
collapsed from 1.00 to 0.27 while an untargeted CLI control held at 1.00,
confirming that the injector degrades a category rather than the system. We report
no detection-lag comparison, because the honest form of the claim is a power
argument that does not need one: sample size scales as $1/\text{effect}^2$, and
a category of share $s$ dilutes its collapse to $s\times$ the effect in the
aggregate, so aggregate-only monitoring needs $1/s^2$ times as many
episodes---$25\times$ at our 20\% share, roughly 4 episodes within the group
against 104 in the aggregate. Aggregate monitoring is not merely slower to see
a minority category fail; at realistic batch sizes it is under-powered to see
it at all. The ratio is scale-free, though the absolute counts depend on a
task-mix-dependent variance estimate. Disaggregated envelopes in turn require
per-group power floors (Section~\ref{sec:trigger}); in one pilot the floor
correctly forced an abstention instead of a confident misattribution.
\subsection{Attribution Quality Across All Cells}
\label{sec:attrquality}

Because every injected cell carries ground truth, the funnel's output is scored
as injected versus attributed mode with abstention as a scored outcome---to our
knowledge the first such accounting for a remediation protocol. Across the four
cycles opened under the frozen mapping table and final contracts, three
produced an attribution and all three were correct (\code{behavioral\_drift} in
$\tau^2$; \code{alignment\_creep} twice in ChatDev, at confidence 0.61 and
0.73); one abstained. Accuracy excluding abstentions is 3/3 at a 25\%
abstention rate. We report the two together, since a stage that resolves every
case by construction can achieve any accuracy demanded of it---and note plainly
that $n=4$ supports an accounting, not a confusion matrix with per-mode recall.
Only APIP runs enter it: B2 does not attribute, and B4 would be diagonal by
fiat.
\subsection{Limitations}

The sharpest limitation is replication: each reported condition is a single
seed on one injection schedule, so we report observed trajectories rather than
rates or confidence intervals, and whether the patch-then-regress dynamic
reproduces is untested.

Every degradation reported here is injected, so the objection that injected
drift merely mirrors its own schedule stands unanswered. Answering it requires
an environment whose world evolves without injection; we discuss the one we
specified for that purpose as future work (Section~\ref{sec:open}).

Injection schedules are scripted, and two contract metrics depend on an LLM
judge; we enforce judge/backbone independence by construction and score one
atomic claim at a time behind a recall-oriented pre-filter, but judge error
remains a bounded unknown. Susceptibility to stale-corpus drift is itself
model-dependent---policy-fact statement rate varied from 3.3\% to 37.5\% across
the backbones we tested---so absolute thresholds do not transfer between
models. That strengthens rather than weakens the argument for behavioral
contracts: thresholds must be calibrated per deployment, which is what the
contract exists to do. Versions are pinned in every run manifest, and all runs,
traces, and the frozen mapping table are released.
% =====================================================================
\section{Open Problems and Future Directions}

\subsection{Multi-Agent Mesh Deployments: The New Hire Problem}
\label{sec:mesh}

The APIP as presented assumes one agent, one contract, one remediation
process. In agent meshes a failure mode emerges that is structurally invisible
to this assumption: \emph{mesh disruption}. A new specialist is introduced and
performs its own task well, yet peer performance declines; the APIP triggers on
the peers, but attribution finds no intrinsic failure mode---no drift, fresh
retrieval, intact alignment. The cause is exogenous, and a single-agent APIP
cannot represent it. Kim et al.\ \cite{kim2025} report performance change
against single-agent baselines from $+80.8\%$ to $-70.0\%$ with a
coordination-saturation effect; the phenomenon is analogous to the
organizational ``new hire problem'' that Event System Theory
\cite{morgeson2015} frames through novelty, disruptiveness, and criticality.

\textbf{We tested this and did not reproduce it.} Inserting a Documentation
Specialist between Programmer and Reviewer executed cleanly---it appeared in
112 of 112 post-injection episodes, produced real documentation, and its output
propagated downstream---yet no peer metric moved: build success 1.00 on both
sides, review engagement 7.64--7.68 before versus 7.63--7.76 after, tool error
rate 0.000. No breach occurred, so no cycle opened and the onboarding arm could
not be evaluated.

We read this as a \emph{precondition} finding rather than a refutation.
Candidate mechanisms include shared-context crowding, role boundary erosion,
distributional shift in shared state, and trust miscalibration. ChatDev builds
each phase's prompt afresh from templates, so there is no bounded context pool
to crowd: crowding is structurally unavailable here, and the mechanism this
environment can test---distributional shift alone---did not degrade peers. Mesh
disruption therefore plausibly requires an architecture with a bounded or
evicting shared context. That is a sharper hypothesis than the one we set out
to confirm, and a concrete requirement on the environment that would test it.
The mesh machinery and the scoped-context onboarding protocol---quarantining a
newcomer's output during a probationary window, analogous to structured
newcomer onboarding \cite{kammeyer2013}---are implemented and released; they
need a substrate where the mechanism can operate, not further code. Whether
attribution stays tractable at hundreds of agents remains open
\cite{krishnan2025}.

\subsection{Open Problems}
\label{sec:open}

\emph{Naturally occurring degradation} is the most immediate next step, since
every result here rests on injected failure. The environment we specified to
remove that dependence---CoffeeBench \cite{coffeebench2026}, whose
agent-operated firms inhabit a persistent economy that evolves on its own---was
not run: its 90-day horizon resists truncation without breaking comparability,
and one replication exceeded our compute budget by an order of magnitude. Its
contract and adapter ship with the harness. It is also where
\code{recalibrated} earns its keep, the protocol there having to separate
degradation warranting intervention from legitimate environmental change.

Three further problems remain open. \emph{Contract specification}: contracts
are architecturally foundational yet absent from most deployments, and
authoring one requires prospective agreement across ML engineering, product,
and legal, with no tooling to support it. \emph{Cause attribution at scale}:
the human adjudication step (Section~\ref{sec:attribution}) does not scale;
automated methods reach only 36.2\% step-level accuracy
\cite{ma2025attribution, zhang2025}, and they attribute individual trajectories
where the APIP needs window-level classification---the gap the funnel of
Section~\ref{sec:funnel} is a first attempt to close. \emph{Organizational
ownership}: in HR the PIP is jointly owned by manager and HR function; in AI
deployments the equivalent is unresolved, and without an ownership model
adoption will stall.

% =====================================================================
\section{Conclusion}

We have introduced the APIP, a post-deployment remediation protocol adapting
the human PIP's lifecycle primitives---trigger, cause attribution, tiered
intervention, bounded window, exit evaluation---with Kirkpatrick's model
separating durable recovery from superficial improvement. Scored end-to-end
against injected ground truth in two live environments, it shows the
distinction is not academic: only durable exit separates remediation from
repeated patching. The harness, artifacts, version pins, and frozen mapping
table are released so the pre-registration claims are checkable. Deployed
agents should be subjects of ongoing accountability rather than static shipped
artifacts---a discipline the field can adopt deliberately, before regulation
makes it mandatory.
% =====================================================================
\begin{thebibliography}{99}
\footnotesize

\bibitem{liang2022}
P. Liang \emph{et al.}, ``Holistic evaluation of language models,'' arXiv
preprint arXiv:2211.09110, 2022.

\bibitem{ganguli2022}
D. Ganguli \emph{et al.}, ``Red teaming language models to reduce harms:
Methods, scaling behaviors, and lessons learned,'' arXiv preprint
arXiv:2209.07858, 2022.

\bibitem{ouyang2022}
L. Ouyang \emph{et al.}, ``Training language models to follow instructions with
human feedback,'' \emph{Advances in Neural Information Processing Systems},
vol.~35, 2022.

\bibitem{pan2025}
M. Z. Pan, N. Arabzadeh, and M. Ellis, ``Measuring agents in production,'' arXiv
preprint arXiv:2512.04123, 2025.

\bibitem{euaiact}
European Parliament and Council, ``Regulation on artificial intelligence (EU AI
Act),'' \emph{Official Journal of the European Union}, 2024.

\bibitem{wachter2017}
S. Wachter, B. Mittelstadt, and C. Russell, ``Counterfactual explanations
without opening the black box: Automated decisions and the GDPR,'' \emph{Harvard
Journal of Law \& Technology}, vol.~31, no.~2, 2017.

\bibitem{kirkpatrick1994}
D. L. Kirkpatrick, \emph{Evaluating Training Programs: The Four Levels}. San
Francisco, CA, USA: Berrett-Koehler, 1994.

\bibitem{cemri2025}
M. Cemri, M. Z. Pan, S. Yang \emph{et al.}, ``Why do multi-agent LLM systems
fail?'' arXiv preprint arXiv:2503.13657, 2025.

\bibitem{sapra2025}
G. Sapra \emph{et al.}, ``AgentCompass: Towards reliable evaluation of agentic
workflows in production,'' arXiv preprint arXiv:2509.14647, 2025.

\bibitem{rath2026}
A. Rath, ``Agent drift: Quantifying behavioral degradation in multi-agent LLM
systems,'' arXiv preprint arXiv:2601.04170, 2026.

\bibitem{zhu2026}
K. Zhu, X. Ye, Z. Han \emph{et al.}, ``AgentDebugX: An open-source toolkit for
failure observability, attribution, and recovery in LLM agents,'' arXiv
preprint arXiv:2607.18754, 2026.

\bibitem{kirkpatrick2016}
J. D. Kirkpatrick and W. K. Kirkpatrick, \emph{Kirkpatrick's Four Levels of
Training Evaluation}. Association for Talent Development, 2016.

\bibitem{kolt2024}
N. Kolt, L. Heim, and M. Anderljung, ``Visibility into AI agents,'' in
\emph{Proc. ACM Conf. Fairness, Accountability, and Transparency (FAccT)}, 2024.

\bibitem{greenblatt2023}
R. Greenblatt, B. Shlegeris, K. Sachan, and F. Roger, ``AI control: Improving
safety despite intentional subversion,'' arXiv preprint arXiv:2312.06942, 2023.

\bibitem{bhatt2025}
A. Bhatt \emph{et al.}, ``Ctrl-Z: Controlling AI agents via resampling,'' arXiv
preprint arXiv:2504.10374, 2025.

\bibitem{schoen2025}
B. Schoen \emph{et al.}, ``Stress-testing control protocols with adaptive
attacks against trusted monitors,'' arXiv preprint arXiv:2510.05367, 2025.

\bibitem{morgeson2015}
F. P. Morgeson, T. R. Mitchell, and D. Liu, ``Event system theory: An
event-oriented approach to the organizational sciences,'' \emph{Academy of
Management Review}, vol.~40, no.~4, pp.~515--537, 2015.

\bibitem{kammeyer2013}
J. D. Kammeyer-Mueller, C. R. Wanberg, A. Rubenstein, and Z. Song, ``Support,
undermining, and newcomer socialization: Fitting in during the first 90 days,''
\emph{Academy of Management Journal}, vol.~56, pp.~1104--1124, 2013.

\bibitem{hendrycks2021}
D. Hendrycks \emph{et al.}, ``Aligning AI with shared human values,'' arXiv
preprint arXiv:2008.02275, 2021.

\bibitem{ma2025attribution}
Y. Ma \emph{et al.}, ``Automatic failure attribution and critical step
prediction for multi-agent systems based on causal inference,'' arXiv preprint
arXiv:2509.08682, 2025.

\bibitem{zhang2025}
S. Zhang \emph{et al.}, ``Which agent causes task failures and when? on
automated failure attribution of LLM multi-agent systems,'' in \emph{Proc. 42nd
Int. Conf. Machine Learning (ICML)}, 2025, arXiv:2505.00212.

\bibitem{agenttrace2026}
``AgentTrace: Causal graph tracing for root cause analysis in multi-agent
systems,'' arXiv preprint arXiv:2603.14688, 2026.

\bibitem{barres2025}
V. Barres, H. Dong, S. Ray, X. Si, and K. Narasimhan, ``$\tau^2$-bench:
Evaluating conversational agents in a dual-control environment,'' arXiv
preprint arXiv:2506.07982, 2025.

\bibitem{qian2024}
C. Qian, W. Liu, H. Liu \emph{et al.}, ``ChatDev: Communicative agents for
software development,'' in \emph{Proc. 62nd Annu. Meeting Assoc. Comput.
Linguistics (ACL)}, 2024.

\bibitem{coffeebench2026}
Sakana AI, ``CoffeeBench: A persistent economic simulation benchmark for LLM
agents,'' \url{https://github.com/SakanaAI/CoffeeBench}, 2026. % TODO: replace
% with the canonical CoffeeBench citation before submission; verify author list
% and venue.

\bibitem{kim2025}
Y. Kim \emph{et al.}, ``Towards a science of scaling agent systems,'' arXiv
preprint arXiv:2512.08296, 2025.

\bibitem{krishnan2025}
N. Krishnan, ``Advancing multi-agent systems through model context protocol,''
arXiv preprint arXiv:2504.21030, 2025.

\end{thebibliography}

\end{document}
